% Generated draft: requires mathematical and source review.
\documentclass[11pt]{article}
\usepackage[margin=25mm]{geometry}
\usepackage{fontspec}
\setmainfont{Latin Modern Roman}
\usepackage{amsmath,amssymb,mathtools,unicode-math}
\setmathfont{Latin Modern Math}
\usepackage{enumitem,needspace,xurl,hyperref}
\setlength{\parindent}{0pt}
\setlength{\parskip}{5pt}
\setlength{\emergencystretch}{3em}
\newcommand{\N}{\mathbb{N}}
\newcommand{\Z}{\mathbb{Z}}
\newcommand{\Q}{\mathbb{Q}}
\newcommand{\R}{\mathbb{R}}
\newcommand{\C}{\mathbb{C}}
\newcommand{\F}{\mathbb{F}}
\newcommand{\A}{\mathbb{A}}
\newcommand{\PP}{\mathbb{P}}
\newcommand{\E}{\mathbb{E}}
\newcommand{\eps}{\varepsilon}
\newcommand{\dd}{\,\mathrm{d}}
\begin{document}
\section*{UnsolvedMath ID: 30004626}
\textbf{Source problem code:} catalogue entry
\subsection{Definitions and mathematical statement}
Let \(X\) be a CAT(0) cube complex with its piecewise Euclidean metric, with no finite upper bound on the dimensions of its cubes, and let a group \(\Gamma\) act on \(X\). For an action by cubical automorphisms, write
\[
\operatorname{Fix}_X(g)=\{x\in X:g\cdot x=x\},\qquad
\operatorname{Fix}_X(\Gamma)=\bigcap_{g\in\Gamma}\operatorname{Fix}_X(g).
\]
An element is \emph{elliptic} if its fixed-point set is nonempty. Under the interpretation that a \emph{locally elliptic action} means that every element is elliptic, the principal question is: under what additional hypotheses on \(\Gamma\), \(X\), and the action does
\[
\bigl(\forall g\in\Gamma\;\exists x_g\in X:\ g\cdot x_g=x_g\bigr)
\quad\Longrightarrow\quad
\bigl(\exists x\in X\;\forall g\in\Gamma:\ g\cdot x=x\bigr)
\]
hold? The supplied statement asks for conditions ensuring this conclusion; it does not specify a single universally quantified conjecture asserting the implication without additional hypotheses.

\textbf{Unresolved qualifications.} The source record does not define ``locally elliptic.'' Another possible convention requires every finitely generated subgroup to have a common fixed point, expressed by
\[
\forall F\subseteq\Gamma\text{ finite},\qquad
\bigcap_{g\in F}\operatorname{Fix}_X(g)\ne\varnothing.
\]
The record does not determine which convention is intended. It also does not specify whether actions must be cubical, whether a fixed point must be a vertex or may be any metric point, or whether completeness, local finiteness, or finite generation of the acting group is required. The displayed principal formulation uses cubical actions and metric fixed points as an explicit interpretation requiring review, rather than as verified source hypotheses.

The original statement contains an additional question omitted from the clean statement: if every element of a finitely generated Cremona subgroup is regularizable, must the subgroup be regularizable? Its mathematical form can be displayed conditionally as follows. For a field \(k\) and a positive integer \(n\), put
\[
\operatorname{Cr}_n(k)=\operatorname{Bir}_k(\mathbb P^n_k).
\]
Fix a permitted class \(\mathcal V\) of \(k\)-varieties birational to \(\mathbb P^n_k\). A subgroup \(H\leq\operatorname{Cr}_n(k)\) is regularizable relative to \(\mathcal V\) if there exist \(Y\in\mathcal V\) and a \(k\)-birational map \(\phi:\mathbb P^n_k\dashrightarrow Y\) such that
\[
\phi H\phi^{-1}\subseteq\operatorname{Aut}_k(Y).
\]
An element \(g\) is regularizable when \(\langle g\rangle\) is regularizable. The additional question asks whether, for finitely generated \(G\leq\operatorname{Cr}_n(k)\),
\[
\bigl(\forall g\in G,\ \langle g\rangle\text{ is regularizable}\bigr)
\quad\Longrightarrow\quad
G\text{ is regularizable}.
\]
The individual regularizations may depend on \(g\); the conclusion requires one common model and conjugacy for all of \(G\). The source does not specify \(k\), \(n\), or the allowed models \(\mathcal V\), including any smoothness or projectivity requirements. This is therefore a conditional formulation of the additional question, not a choice of those missing hypotheses. No equivalence between the two questions is asserted.
\subsection{Short English statement}
The main question asks which hypotheses force an infinite-dimensional CAT(0) cube complex action to have one point fixed by the whole group when the action is locally elliptic. The meaning of local ellipticity and several action hypotheses remain unspecified. The original record also asks whether elementwise regularizability of a finitely generated Cremona subgroup implies simultaneous regularizability; its field, rank, and permitted models require clarification.
\subsection{Sources}
\begin{itemize}
\item Local UnsolvedMath registry, supplied complete source record, permanent ID \(30004626\), exact source problem code \texttt{catalogue entry}, in the Oberwolfach Reports Open Problems collection. The supplied proposed year is 2021; no individual proposer is named.
\item Supplied report citation: \emph{Nonpositively Curved Complexes (2022)}, Oberwolfach Reports, DOI \texttt{10.4171/owr/2021/8}, pp.~383--417. \url{https://doi.org/10.4171/owr/2021/8}. The record attributes the problem to workshop open-problem discussions collected in this report. Its title year and DOI year are retained as supplied.
\item Lonjou--Przytycki--Urech, ``Cremona subgroups'' (supplied reference label). \url{https://doi.org/10.5802/jep.271}.
\item Genevois--Lonjou--Urech, ``finite-field regularisation'' (supplied reference label). \url{https://arxiv.org/abs/2110.14605}.
\item Supplied literature review, checked exactly \texttt{2026-08-22}: the record labels the problem open and states that the regularizability implication for finitely generated Cremona subgroups and the locally elliptic fixed-point question remain unresolved over general fields or infinite-dimensional complexes. This is supplied dated curation derived from the OWR material, not an independently verified literature assessment. The listed references and their applicability require independent checking.
\item The record labels the statement \texttt{corrected\_verified} and reports a check against the official Oberwolfach report preserving the clean statement's hypotheses, notation, and target. This is supplied verification metadata; no source visit or independent verification is claimed here.
\end{itemize}
\end{document}
